\documentclass[a4paper, 12pt]{mksheetphhd}

\title{Animationen}
\author{Marvin Ködding}
\usepackage{minted}
\usepackage{multicol}

\everymath{\displaystyle}

\begin{document}
	\printtitle
	
	\vspace*{.5cm}
	
	\begin{aufgabe}
		Zeichne zwei Kreise, die links und rechts auf dem Bildschirm liegen. Verbinde ihre jeweiligen rechten und linken kritischen Punkte mit Linien. Verwende \verb|get_critical_point| und Line.
	\end{aufgabe}
	
	\begin{aufgabe}
		Zeige einen längeren Satz als Text-Objekt. Lasse ein Wort darin mit \verb|FocusOn| hervorheben und anschließend mit \verb|Circumscribe| umranden.
	\end{aufgabe}
	
	\begin{aufgabe}
		Erstelle das Wort \texttt{VISUALISIERUNG}, wobei jeder Buchstabe eine andere Farbe bekommt und mit AddTextLetterByLetter einzeln erscheint. Danach lasse die Buchstaben mit FadeOut verschwinden – aber rückwärts Buchstabe für Buchstabe (also letzter Buchstabe zuerst).
		
		{\it Hinweis: Schleife über die Buchstaben oder Liste der Buchstaben in umgekehrter Reihenfolge (\verb|list[::-1]| gibt eine Liste in umgekehrter Reihenfolge aus)}
	\end{aufgabe}
	
	\begin{aufgabe}
		Erstelle ein 5×5-Raster aus Kreisen \verb|(VGroup(...).arrange_in_grid())| und lasse einen zufälligen Kreis kurz mit Flash aufleuchten. Nutze \verb|random.choice(...)|.
	\end{aufgabe}
	
	\begin{aufgabe}
		Erstell eine Szene, in der die Produktregel der Differentialrechnung schrittweise aufgebaut und animiert wird.
		
		Folgende Schritte sollen umgesetzt werden:
		\begin{enumerate}
			\item Starte mit der Formel \(\frac{\mathrm{d}}{\mathrm{d}x} \bigl[f(x) \cdot g(x)\bigr]\).
			\item Ersetze sie schrittweise durch \(\frac{\mathrm{d}}{\mathrm{d}x}f(x) \cdot g(x) + f(x) \cdot \frac{\mathrm{d}}{\mathrm{d}x}g(x)\) oder nutze eine andere Animation.
			\item Wenn du magst, kannst du durch \verb*|Indicate| o.Ä. die zu ersetzenden Parts kennzeichnen.
		\end{enumerate}
	\end{aufgabe}
	
	\newpage
	
	Beispiellösung:
	\begin{minted}[bgcolor=gray!10, breaklines]{python}
class ProductRuleAnimation(Scene):
  def construct(self):
    # Schritt 1: Ursprungsformel
    start = MathTex(r"{{\frac{d}{dx}}} \left[ {{f(x)}} \cdot {{g(x)}} \right]", font_size=72).shift(UP)
    self.play(Write(start))
    self.add(index_labels(start))
    self.wait()
    print(len(start))
    
    # Schritt 2: Produktregel
    result = MathTex(
    r"{{\frac{d}{dx}}} {{f(x)}} \cdot {{g(x)}} + {{f(x)}} \cdot {{\frac{d}{dx}}}{{g(x)}}",
    font_size=72
    ).shift(DOWN)
    self.add(index_labels(result))
    
    self.play(
    	TransformFromCopy(start[0], result[0]),
    	TransformFromCopy(start[2], result[2]),
    )
    self.play(FadeIn(result[3]))
    self.play(TransformFromCopy(start[4], result[4]))
    self.play(FadeIn(result[5]))
    self.play(TransformFromCopy(start[2], result[6]))
    self.play(FadeIn(result[7]))
    self.play(
    	TransformFromCopy(start[0], result[8]), TransformFromCopy(start[4], result[10]))
    self.wait()
	\end{minted}
	
	\printlicense
	
	\printsocials
	
\end{document}